\section{Conclusion}
\label{sec:conclusion}

We asked what unit of evidence a temporal knowledge graph forecaster should read, and measured the answer before modelling it. The dyad, the ordered entity pair with the typed stream of every event between its members, carries information that the query pair does not: a quarter of ICEWS queries have an answer that only the dyad can see, and a counter that reads the dyad with no training matches a published model that reads the query pair with a great deal of it. \model{} turns that measurement into a model: a superposition of intensities, with a Transformer over the dyad's event stream at its centre, a path intensity that composes a stream with a typed edge, and attention across the candidates of a query. Over three seeds it sets the best published MRR and \hone{} on YAGO and WIKI, the best \hone{} on GDELT, and exceeds every method with released code on ICEWS18 and ICEWS14s. The stratified analysis shows each intensity acting where the data said it would, and shows as plainly where the model does not act: a third of ICEWS queries whose answer never interacted with the subject remain largely out of reach, and three recent methods without released code report higher aggregate numbers on those benchmarks. Reaching the cold strata, and settling the comparison with those methods under one protocol, are the two open ends of this work. The code, the diagnostic and every result file are released so that both can be taken up.
